\documentclass[12pt]{article}
\usepackage[utf8]{inputenc}
\usepackage[T1]{fontenc}
\usepackage{hyperref}
\title{Tigr\'e Language Resources for K-12 Curriculum Development}
\author{Compiled by Claude Sonnet 5 (claude-sonnet-5) for Andr\'as Kornai}
\date{}
\begin{document}
\maketitle
\section{Introduction}
Tigr\'e (ISO 639-3 tig) is an Afro-Asiatic (Ethiosemitic) language of Eritrea and eastern Sudan, written in the Ge'ez script. Speaker figures vary: about 800,000 in one report, and about 35\% of Eritrea's population in another source. The search surfaced policy and background material but no actual Tigr\'e textbook titles.
\section{Primary Educational Resources}
No named textbooks were found. A Leeds article says Eritrea's mother-tongue policy allows elementary schooling in each of nine languages, including Tigr\'e. It also says textbook production continued in all nine languages after 1991: \url{https://lucas.leeds.ac.uk/keyword/literacy/}
\section{Linguistic and Reference Materials}
\begin{itemize}
\item Endangered Languages entry for Tigr\'e: \url{https://www.endangeredlanguages.com/lang/4603}
\item Rhode Island College language project report on Tigre: \url{https://ritell.org/Resources/Documents/language\%20project/Tigre.pdf}
\end{itemize}
\section{Religious and Community Resources}
None found.
\section{Digital and Audio Resources}
None found.
\section{Curriculum Development Context}
The policy framework supports Tigr\'e as a language of instruction in elementary school. The language has a written history of more than 100 years. An opinion piece by a PEN Eritrea writer argues that Tigr\'e speakers increasingly choose Arabic in education and public life: \url{https://www.upplittmagasin.se/artikel/tigre---a-language-betrayed-from-within}. Further work should look for Eritrean Ministry of Education and Eritrean Institute for Educational Research materials.
\end{document}